\section{Complete gauge fixing and the Fredholm obstruction}\label{sec:fredholm}
We continue with the common normalization and reduced data above.
Write $\mathcal G^s=\Ran(K_\Sigma:\mathcal H_1^{s+1}\to\mathcal H_2^s)$ and $\mathcal C^s=\Ker K_\Sigma^\dagger$; omission of the superscript denotes smooth data.
Set $Y^s=H^s(\Sigma;E_1)\oplus H^s(\Sigma;E_1)$.
The readout $R:\mathcal H_2^s\to Y^s$ is surjective: it is the composition of the graph-coordinate isomorphism with $l_N\oplus l_N$, and $l_N$ has a smooth fibrewise right inverse.
For example, its prescribed trace can be supplied by a multiple of $g_0$ and its mixed component can be specified independently.
Surjectivity is an exact statement, including low frequencies.

\subsection{The exact boundary matrix}
Let $r=\tfrac12h_0^{-1}\dot h_0$ and $H=\tr r$.
For reduced one-form data $w=(w_t,w_\Sigma)$ the original covector is $a(t)W(t)w(t)$, where $a'/a=-H/2$ and spatial parallel transport satisfies $W'=rW$ at the reference slice.
The Gaussian Christoffel symbols are $\Gamma^j_{ti}=r_i{}^j$ and $\Gamma^t_{ij}=r_{ij}$.
Direct substitution gives
\begin{align*}
\tr_g(I d_s\xi)|_\Sigma
 &=\partial_t w_t+\tfrac H2w_t-\operatorname{div}_\Sigma w_\Sigma,\\
2(d_s\xi)_{t\Sigma}|_\Sigma
 &=\partial_t w_\Sigma-\tfrac H2w_\Sigma-rw_\Sigma+d_\Sigma w_t.
\end{align*}
Consequently, with $\delta_\Sigma=-\operatorname{div}_\Sigma$ and adjoints taken in the positive Riemannian product,
\begin{equation}\label{eq:B}
B=\begin{pmatrix}H/2&\delta_\Sigma\\d_\Sigma&-H/2-r\end{pmatrix},\qquad B^*=B,
\end{equation}
and
\begin{equation}\label{eq:readoutK}
l_N\pi_0K_\Sigma c_1^\pm
=\tfrac1{\sqrt2}(i b_1^\pm+B)\pi_0c_1^\pm.
\end{equation}
The coefficient of $r$ in~\eqref{eq:B} follows from the displayed differentiation.
In the primary rendering~\cite[Equations (4.10)--(4.11)]{Gerard2024Compact}, the derivative formula agrees with this coefficient whereas the following matrix displays $-r/2$.
We use~\eqref{eq:B}, with $r=\tfrac12h^{-1}\dot h$ fixed as above.

The scalar row of $K$ is removed exactly by $l_N$.
Using the smoothing intertwining defect and writing $L_1f=(\pi_0c_1^+f,\pi_0c_1^-f)$, we obtain
\begin{equation}\label{eq:RK}
A:=RK_\Sigma
=\tfrac1{\sqrt2}\begin{pmatrix}A^+&0\\0&A^-\end{pmatrix}L_1+S,
\quad A^\pm=i b_1^\pm+B,\quad S\in\Psi^{-\infty}.
\end{equation}
Thus all coupled lower-order dependence is included in $b_1^\pm$ and $S$; no unaccounted order-zero error is discarded.

\begin{proposition}\label{prop:fredholm}
For every $s\in\mathbb R$, $A:\mathcal H_1^{s+1}\to Y^s$ is Fredholm of index zero.
Its kernel and distributional cokernel are finite dimensional, smooth and independent of $s$.
The gauge range $\mathcal G^s$ is closed.
The induced charge on $\mathcal C/\mathcal G$ is nondegenerate.
\end{proposition}
\begin{proof}
Each $A^\pm:H^s\to H^{s-1}$ is elliptic: its principal symbol is $\pm i|\eta|+\sigma_1(B)$ and $\sigma_1(B)$ is Hermitian.
At $s=1/2$, the imaginary part of $(u,A^\pm u)$ is $(u,\operatorname{Re}b_1^\pm u)$, with a fixed sign and a coercive $H^{1/2}$ bound.
The forms rotated by $\mp i$ are coercive bounded forms on $H^{1/2}$.
Lax--Milgram gives injectivity and surjectivity $H^{1/2}\to H^{-1/2}$.
The adjoint estimate has the opposite imaginary sign and excludes a cokernel as well.
Elliptic regularity gives the isomorphism for every $s$.
Because $L_1:\mathcal H_1^{s+1}\to H^{s+1}\oplus H^{s+1}$ is an isomorphism, the diagonal term in~\eqref{eq:RK} is invertible onto $Y^s$.
Its smoothing perturbation is Fredholm of index zero, with smooth kernel and cokernel.
The perturbation can create both, so invertibility of $A$ is not assumed.

Choose a closed complement $X_0$ to $\Ker A$.
The restriction $A:X_0\to\Ran A$ is a bounded isomorphism.
If $K_\Sigma x_n$ converges for $x_n\in X_0$, then $Ax_n=RK_\Sigma x_n$ converges, hence $x_n$ converges.
Thus $K_\Sigma X_0$ is closed; adding the finite-dimensional smooth space $K_\Sigma\Ker A$ proves closedness of $\mathcal G^s$.

At $s=1/2$, the charge is a strong nondegenerate form: its off-diagonal entries identify $H^{1/2}$ and $H^{-1/2}$ with the appropriate dual spaces.
The charged-adjoint identity gives $\mathcal C^{1/2}=(\mathcal G^{1/2})^{q_J}$.
Closedness then implies $(\mathcal C^{1/2})^{q_J}=\mathcal G^{1/2}$, proving nondegeneracy of the quotient.
This also gives the assertion for smooth data.
Indeed, after Sobolev conjugation, the principal symbol of $K_\Sigma$ is injective because that of $RK_\Sigma$ is invertible.
Its positive normal operator is elliptic; its generalized inverse and range projection preserve smooth sections.
The corresponding constraint projection, obtained using the charge isomorphism, also preserves smooth sections.
Smooth constrained data are therefore dense in $\mathcal C^{1/2}$, and a smooth gauge-range element has a smooth gauge preimage.
\end{proof}

\section{Positive covariances, including the singular case}\label{sec:covariances}
We now construct a bounded gauge correction on each graded Sobolev space.
Here $Q=q_J$ denotes the physical Cauchy charge as a matrix in the positive auxiliary pairing.
For a field-data operator $T$, put $T^\dagger=Q^{-1}T^*Q$; for a readout map into the positive target, put $R^\dagger=Q^{-1}R^*$.
The distinction from $K_\Sigma^\dagger=q_1^{-1}K_\Sigma^*Q$ is essential.

\subsection{The regular case}
If $A$ is invertible, define
\[
T=1-K_\Sigma A^{-1}R.
\]
Then $RT=0$, $T^2=T$, $TK_\Sigma=0$, and $T$ preserves $\mathcal C$ because $K_\Sigma^\dagger K_\Sigma=0$.
For $f\in\mathcal C$, $Tf-f\in\mathcal G$.
The positive slice proved above gives positive Hermitian covariances
\[
\lambda_\Sigma^\pm=\pm T^*Q c_2^\pm T.
\]
They annihilate $\mathcal G$ in both arguments, and their difference is $Q$ on $\mathcal C\times\mathcal C$ by charge orthogonality to $\mathcal G$.
Moreover $c_2^\pm Tc_2^\mp$ is smoothing: $R$ separates the target frequencies, $A^{-1}$ is diagonal modulo smoothing in graph coordinates, and $K_\Sigma$ intertwines the frequency splittings modulo smoothing.

\subsection{The full kernel and cokernel construction}
Take a positive auxiliary $L^2$ basis $u_1,\ldots,u_n$ of $\Ker A^*$ and set
\[
v_j=R^\dagger u_j=Q^{-1}R^*u_j.
\]
Elliptic regularity makes all these vectors smooth, and surjectivity of $R$ makes the $v_j$ independent.
The exact charged-adjoint identity
\[
0=A^*u_j=K_\Sigma^*R^*u_j=q_1K_\Sigma^\dagger v_j
\]
places every $v_j$ in $\mathcal C$.
There is no assertion that $R^\dagger u_j$ has a vanishing scalar component.
The solvable-readout subspace is
\[
\mathcal H_{\mathrm{reg}}^s
=\{f:Rf\in\Ran A\}
=\{f:q_J(v_j,f)=0\text{ for every }j\}.
\]

Choose the basis so that $v_1,\ldots,v_p$ represent a complement of $V\cap\mathcal G$ in $V=\operatorname{span}(v_j)$ and $v_{p+1},\ldots,v_n\in\mathcal G$.
Nondegeneracy of $\mathcal C/\mathcal G$ and smooth density supply smooth $w_1,\ldots,w_p\in\mathcal C$ with an invertible pairing matrix against the first $p$ vectors.
Their pairings against the remaining vectors vanish.
Ambient charge nondegeneracy completes them by smooth $w_{p+1},\ldots,w_n$ to a full invertible matrix.
The block triangular inverse normalizes the dual family while keeping the first $p$ vectors in $\mathcal C$:
\begin{equation}\label{eq:dual}
q_J(v_j,w_i)=\delta_{ji},\qquad
w_i\in\mathcal C\ (i\le p),\quad v_j\in\mathcal C\ (j\le n).
\end{equation}
All dual pairings in~\eqref{eq:dual} use the physical charge.

Define smooth finite-rank projections
\[
\pi f=\sum_{j=1}^n w_jq_J(v_j,f),\qquad
\widetilde\pi f=\sum_{j=1}^p w_jq_J(v_j,f).
\]
Then $\Ker\pi=\mathcal H_{\mathrm{reg}}$, $\pi K_\Sigma=0$, and $\pi$ preserves $\mathcal C$.
On $\mathcal C$, $\pi=\widetilde\pi$, since the omitted $v_j$ belong to $\mathcal G$.
In addition,
\begin{equation}\label{eq:piannihilation}
\widetilde\pi K_\Sigma=\widetilde\pi^\dagger K_\Sigma=0.
\end{equation}
The second identity uses $w_j\in\mathcal C$ for $j\le p$.
The full $\pi^\dagger$ need only preserve $\mathcal G$ and is not asserted to annihilate it.

Let $\pi_1$ be a smooth finite-rank projection onto $\Ker A$ and let $G_A$ be the generalized inverse with range $\Ker\pi_1$, extended to all target data through a smooth cokernel projection.
Thus $G_AA=1-\pi_1$.
Choose a smooth finite-rank projection $\pi_2$ onto $K_\Sigma\Ker A=\mathcal G\cap\Ker R$.
Set
\begin{equation}\label{eq:Tsingular}
T_{\mathrm{reg}}=1-K_\Sigma G_AR\quad\hbox{on }\mathcal H_{\mathrm{reg}},
\qquad T=(1-\pi_2)T_{\mathrm{reg}}(1-\pi).
\end{equation}
These maps have the required boundedness on every $\mathcal H_2^s$.
Direct calculation gives $RT=0$, and $T$ is a projection onto $\Ker R\cap\Ker\pi_2$.
It preserves $\mathcal C$.
Since $T_{\mathrm{reg}}K_\Sigma=K_\Sigma\pi_1$, one has
\[
TK_\Sigma=(1-\pi_2)K_\Sigma\pi_1=0.
\]
The $\pi_2$ factor removes residual gauge fields from $\Ker A$ and cannot be omitted.
For constrained data,
\begin{equation}\label{eq:Tquotient}
Tf=(1-\pi)f\pmod{\mathcal G}
=(1-\widetilde\pi)f\pmod{\mathcal G}.
\end{equation}
All extra projections are smooth finite rank, so $c_2^\pm Tc_2^\mp$ remains smoothing.

Define the finite-rank Hermitian smoothing operator
\begin{equation}\label{eq:nu}
\nu=Q-(1-\widetilde\pi)^*Q(1-\widetilde\pi)
=Q\widetilde\pi+\widetilde\pi^*Q-\widetilde\pi^*Q\widetilde\pi.
\end{equation}
Equations~\eqref{eq:piannihilation} imply $\nu K_\Sigma=0$.
The positive and negative spectral parts in the auxiliary positive $L^2$ space,
\[
\nu_+=\tfrac12(|\nu|+\nu),\qquad
\nu_-=\tfrac12(|\nu|-\nu),
\]
are positive smooth finite-rank forms with $\nu=\nu_+-\nu_-$.
Their nonzero eigenvectors lie in the smooth range of $\nu$.
Every gauge vector belongs to its zero eigenspace, so $\nu_\pm$ both annihilate $\mathcal G$.
This conclusion follows by spectral calculus and requires no sign assumption on the finite physical obstruction space.

\begin{proposition}\label{prop:covariances}
The covariances
\begin{equation}\label{eq:lambdas}
\lambda_\Sigma^\pm=\pm T^*Q c_2^\pm T+\nu_\pm
\end{equation}
are positive and Hermitian, annihilate $\mathcal G$ and satisfy
$\lambda_\Sigma^+-\lambda_\Sigma^-=Q$ on $\mathcal C\times\mathcal C$.
Their cross-frequency errors are smoothing.
\end{proposition}
\begin{proof}
The first term in~\eqref{eq:lambdas} is positive because $\Ran T\subset\Ker R$ and the exact slice identity applies.
Hermiticity follows from physical-charge self-adjointness of $c_2^\pm$.
Gauge annihilation follows from $TK_\Sigma=0$ and the spectral argument.
Finally~\eqref{eq:Tquotient}, charge orthogonality to $\mathcal G$, and~\eqref{eq:nu} give, on constrained data,
\[
\lambda_\Sigma^+-\lambda_\Sigma^-
=T^*QT+\nu
=(1-\widetilde\pi)^*Q(1-\widetilde\pi)+\nu=Q.
\]
The finite-rank terms are smoothing and the remaining frequency assertion was proved after~\eqref{eq:Tsingular}.
\end{proof}

The construction treats gauge stabilizers $K_\Sigma\xi=0$, nonzero residual gauge fields $K_\Sigma\Ker A$, and genuine physical obstruction classes $V/(V\cap\mathcal G)$.
It follows the singular method of~\cite[Theorem 5.2]{Gerard2024Compact}, with a single physical charge throughout the dual matrix and defect~\eqref{eq:nu}.
Changing to the untrace-reversed charge in either step would not establish the stated identities.

\section{Spacetime lift and real physical state}\label{sec:lift}
Let $G_2=G_{2,\mathrm{ret}}-G_{2,\mathrm{adv}}$, with the charge convention fixed consistently with the Green identity, and let $\rho_2$ be the reduced trace.
For $f,h\in C_c^\infty(M;E_2)$ define
\begin{equation}\label{eq:lift}
\Lambda^\pm(f,h)=\lambda_\Sigma^\pm(\rho_2G_2f,\rho_2G_2h).
\end{equation}
The Green operator maps compact tests continuously to smooth spacelike-compact solutions, with advanced and retarded support in the corresponding causal cones.
Their Cauchy traces are smooth; compactness of $\Sigma$ admits every smooth datum.
Only a Gaussian collar is used for the microlocal construction, and neither a global Gaussian foliation nor geodesic completeness is needed.
Smooth finite-rank Cauchy terms lift to smooth kernels, since evolution sends smooth data to smooth global solutions.

\subsection{The physical test quotient}
Let $\Ker_cK^\dagger$ denote compactly supported smooth tests satisfying the constraint.
Intertwining implies that $\rho_2G_2$ maps this space into $\mathcal C$.
Furthermore $G_2Ph=-KG_1K^\dagger h$, so $P C_c^\infty(M;E_2)$ maps into the gauge range.
Hence~\eqref{eq:lift} descends to
\begin{equation}\label{eq:testquotient}
\frac{\Ker_cK^\dagger}{P C_c^\infty(M;E_2)}
\simeq\frac{\mathcal C}{\mathcal G}.
\end{equation}
We give the support-sensitive proof of the isomorphism.

Choose a smooth temporal cutoff $\chi$ whose transition lies in a compact Cauchy slab, with constant values zero in the past and one in the future.
Reverse its sign if required by the propagator convention.
Given $D_2u=0$ and $K^\dagger u=0$, put
\[
f=D_2(\chi u),\qquad v=K^\dagger(\chi u).
\]
Both have compact support because the equations vanish outside the transition.
Then $K^\dagger f=D_1v$, and
\[
f_{\mathrm{phys}}=f-Kv\in\Ker_cK^\dagger,
\qquad G_2f_{\mathrm{phys}}=u-KG_1v.
\]
The cutoff Green identity proves surjectivity onto the physical quotient, without requiring surjectivity onto each fixed constrained representative.

For injectivity, suppose $f\in\Ker_cK^\dagger$ and $G_2f=Kw$ with $D_1w=0$.
Set $\eta=D_1(\chi w)$ so that $G_1\eta=w$.
Then $G_2(f-K\eta)=0$.
Exactness of the normally hyperbolic causal sequence supplies compact $h$ satisfying $f-K\eta=D_2h$.
Applying $K^\dagger$ yields
\[
D_1(\eta+K^\dagger h)=0.
\]
This homogeneous solution has compact spacetime support and therefore vanishes by Cauchy uniqueness.
Consequently $f=D_2h-KK^\dagger h=Ph$.
Conversely every such $Ph$ maps to a gauge solution as above.
This proves~\eqref{eq:testquotient} using only causal-sequence exactness, cutoff Green identities and uniqueness for normally hyperbolic equations.

\subsection{Microlocal regularity and the neutral real field}
Because $c_2^\pm$ is $Q$-self-adjoint and $c_2^\pm Tc_2^\mp$ is smoothing, the same cross-frequency assertion holds for $T^\dagger$.
Thus
\[
U_2T^\dagger c_2^\pm T
=U_2c_2^\pm T^\dagger T\quad\pmod{\text{smooth kernel}}.
\]
The factor on the right acting on Cauchy data is graded pseudodifferential.
The microlocal Cauchy-lift criterion~\cite[Proposition 11.1.1]{Gerard2019Microlocal}, used in~\cite[proof of Theorem 5.1]{Gerard2024Compact}, together with propagation of singularities, gives
\[
\WF'(\Lambda^\pm)\subset\mathcal C_{\mathrm{Had}}^\pm
\subset\mathcal N^\pm\times\mathcal N^\pm.
\]
Here $\mathcal C_{\mathrm{Had}}^\pm$ is the relation of cotangent vectors on the same null bicharacteristic with the indicated energy sign, and $\WF'$ reverses the second covector.
Spatial pseudodifferential kernels are conormal to the diagonal and evolution propagates that relation along null bicharacteristics.
The equation removes any noncharacteristic contribution in the local parametrix statement.
The finite-rank corrections introduce no new singularities.
We assert this Hadamard inclusion for the redundant field representatives and exact commutation relations on physical tests; no equality on every component of the redundant bundle is required.

Arbitrary smoothing choices need not preserve the real-field structure.
On reduced Cauchy data let $C(f_0,f_1)=(\overline f_0,-\overline f_1)$ and define $a^C(f,h)=\overline{a(Cf,Ch)}$.
The real differential gauge complex is preserved and $Q^C=-Q$.
Replace the covariances by
\[
\lambda_{\mathrm{real}}^\pm
=\tfrac12\bigl(\lambda_\Sigma^\pm+(\lambda_\Sigma^\mp)^C\bigr).
\]
This preserves positivity, gauge annihilation and the exact physical commutation relations.
Conjugation exchanges the two characteristic signs, so the Hadamard property is preserved.
The pair now satisfies $(\lambda_{\mathrm{real}}^+)^C=\lambda_{\mathrm{real}}^-$ and defines a quasifree state on the neutral real physical CCR algebra.
Explicitly, on the $C$-fixed real quotient set $\sigma=-iQ$ and $\mu=\tfrac12(\lambda_{\mathrm{real}}^++\lambda_{\mathrm{real}}^-)$.
The relation $Q^C=-Q$ makes $\sigma$ real and antisymmetric, while positivity of $\lambda_{\mathrm{real}}^\pm=\mu\pm i\sigma/2$ gives the quasifree uncertainty inequality.
With the Weyl convention $W(f)W(h)=e^{-i\sigma(f,h)/2}W(f+h)$, the state is $\omega(W(f))=e^{-\mu(f,f)/2}$.
Together with Proposition~\ref{prop:covariances} and~\eqref{eq:testquotient}, this proves Theorem~\ref{thm:main}.
